Hacer este análisis nos sirvió para comprobar que la notación Big-O realmente ayuda a predecir cómo se van a comportar los algoritmos cuando procesan millones de datos. Pero la gran conclusión es que la teoría no lo es todo. En la práctica, las pequeñas decisiones al momento de programar (como usar la partición de 3 vías para evitar los repetidos o frenar la recursividad en Strassen) son las que realmente salvan el código y hacen que los algoritmos funcionen rápido en la vida real.